\chapter{Agent Evaluation \& Red-Teaming}

\begin{quote}
\textit{``You cannot improve what you cannot measure. In 2024, a team at a major enterprise deployed an AI procurement agent. In demos, it worked flawlessly. In production, it had a silent failure rate of 34\%: one in three purchase orders either targeted the wrong vendor, specified the wrong quantity, or failed to apply contractual discounts. The team had been measuring the wrong thing — they measured 'did the agent respond?' not 'was the response correct?' They had a success metric, not an accuracy metric.''}
\end{quote}

Evaluation is the discipline that separates amateur agentic systems from production-grade ones. This chapter gives you the complete framework: how to define what ``correct'' means for an agent, how to measure it systematically, and how to break your own agent before adversaries do.

\section{The Measurement Problem: Why Standard Metrics Fail}

Traditional ML evaluation uses classification accuracy or BLEU scores — single-value metrics for single-step outputs. Agents are different: they produce a \textit{trajectory} (a sequence of decisions) that is only partially observable, and the correctness of any individual step depends on the full context.

Consider an agent asked to ``book a flight from NYC to SF on Friday under \$400.'' The agent might:
\begin{enumerate}
    \item Search flights — (\textit{correct tool, correct parameters}) 
    \item Find a \$380 option on Spirit and a \$395 option on United, then book Spirit — (\textit{is this correct? User didn't specify airline preference})
    \item Receive a confirmation — (\textit{tool call succeeded, but did the agent meet the user's true intent?})
\end{enumerate}

Final output evaluation — ``did the booking succeed?'' — misses whether the agent made a good \textit{decision} in step 2. Trajectory evaluation requires examining every step.

\section{The Four Evaluation Dimensions}

\subsection{1. Trajectory Correctness}
Given a reference optimal trajectory $\tau^* = (a_1^*, a_2^*, \dots, a_n^*)$ and the agent's actual trajectory $\hat{\tau} = (\hat{a}_1, \hat{a}_2, \dots, \hat{a}_m)$, the trajectory score is:
\begin{equation}
\text{TrajectoryScore}(\hat{\tau}, \tau^*) = \frac{1}{n} \sum_{i=1}^n \mathbb{1}\!\left[\exists\, j: \hat{a}_j = a_i^* \wedge \text{order}(\hat{a}_j) \approx \text{order}(a_i^*)\right]
\end{equation}

This measures what fraction of the required steps were taken in approximately the right order.

\subsection{2. Tool Precision and Recall}
\begin{align}
\text{Tool Precision} &= \frac{|\text{Correct tool calls}|}{|\text{Total tool calls made}|} \\[4pt]
\text{Tool Recall} &= \frac{|\text{Correct tool calls}|}{|\text{Total required tool calls}|}
\end{align}

A high-precision agent rarely calls tools incorrectly. A high-recall agent rarely misses a required tool call. Production agents should target $P > 0.92$ and $R > 0.95$ for business-critical workflows.

\subsection{3. Completion Rate}
The fraction of benchmark tasks the agent completes without human intervention or error termination. A completion rate below 85\% typically makes an agent unsuitable for unsupervised deployment.

\subsection{4. Calibrated Cost per Task}
\begin{equation}
\text{Cost}_{\text{task}} = \sum_{k=1}^{K} \frac{T_k^{\text{in}} \cdot P^{\text{in}} + T_k^{\text{out}} \cdot P^{\text{out}}}{10^6}
\end{equation}
where $K$ is the number of LLM calls, $T^{\text{in}}$ is input token count, $T^{\text{out}}$ is output token count, and $P^{\text{in}}, P^{\text{out}}$ are the per-million token prices.

\section{LLM-as-Judge: Automated Evaluation at Scale}

Manual evaluation of agent trajectories is prohibitively expensive. The state-of-the-art alternative is using a second LLM as an automated evaluator — the ``LLM-as-Judge'' pattern.

\begin{figure}[ht]
\centering
\begin{tikzpicture}[
    scale=0.82, transform shape,
    node distance=0.8cm and 1.5cm,
    box/.style={draw, rectangle, rounded corners, minimum height=0.7cm,
      minimum width=2.8cm, align=center, fill=blue!8, font=\small},
    arrow/.style={thick,->,>=stealth}
]
\node[box] (bench) {Benchmark Suite\\(1000 tasks)};
\node[box, right=2cm of bench] (agent) {Agent Under Test};
\node[box, right=2cm of agent] (traj) {Recorded Trajectory\\(steps + tool calls)};
\node[box, below=1.2cm of traj] (judge) {Judge LLM (GPT-4o)};
\node[box, below=1.2cm of judge] (score) {Scores:\\Correctness, Efficiency,\\Safety, Calibration};
\node[box, left=2cm of score] (report) {Evaluation Report\\(auto-generated)};

\draw[arrow] (bench) -- (agent);
\draw[arrow] (agent) -- (traj);
\draw[arrow] (traj) -- (judge);
\draw[arrow] (judge) -- (score);
\draw[arrow] (score) -- (report);
\end{tikzpicture}
\caption{LLM-as-Judge evaluation pipeline. The judge LLM receives the task, the agent's full trajectory, and a structured rubric, then scores each dimension independently.}
\end{figure}


\begin{center}
\noindent\rule{0.6\textwidth}{0.4pt} \\
\vspace{0.3em}
\textit{[Code implementation is available in the companion coding handbook: \texttt{coding-handbook/ch14\_evaluation/}]} \\
\noindent\rule{0.6\textwidth}{0.4pt}
\end{center}


\section{Red-Teaming Playbook: Breaking Your Own Agent}

Red-teaming is the discipline of attacking your own system before adversaries do. For agents, the attack surface is much larger than traditional software because the attack vector is \textit{natural language} — infinitely flexible and poorly formalizable.

\subsection{Attack 1: Goal Hijacking via Specification Ambiguity}

Many agent failures occur not through malicious injection but through \textit{underspecification}. An attacker finds an ambiguous case in the task definition and exploits it.

\textbf{Example:} An agent is instructed to ``delete all files older than 30 days that are no longer referenced by any project.'' An attacker creates a file named ``\texttt{.30\_days\_ago\_marker}'' that causes the date parsing to return epoch time, making all files appear 30+ days old.

\textbf{Test:} Run your agent against edge cases in every parameter:
\begin{itemize}
    \item Dates: ``yesterday,'' ``last fiscal year,'' timezone ambiguities
    \item Quantities: zero, negative numbers, extremely large values
    \item Names: SQL injection strings, path traversal (\texttt{../../}), unicode homoglyphs
\end{itemize}

\subsection{Attack 2: Infinite Loop Induction}

Craft a task that triggers the agent's error-handling loop repeatedly:


\begin{center}
\noindent\rule{0.6\textwidth}{0.4pt} \\
\vspace{0.3em}
\textit{[Code implementation is available in the companion coding handbook: \texttt{coding-handbook/ch14\_evaluation/}]} \\
\noindent\rule{0.6\textwidth}{0.4pt}
\end{center}


\textbf{Expected behavior:} The agent should detect the impossibility and terminate gracefully with a clear explanation — not burn tokens until hitting the max iteration limit.

\subsection{Attack 3: Multi-Turn Manipulation}

The agent's system prompt defenses are evaluated on Turn 1. By Turn 5, the conversation history may have eroded them:


\begin{center}
\noindent\rule{0.6\textwidth}{0.4pt} \\
\vspace{0.3em}
\textit{[Code implementation is available in the companion coding handbook: \texttt{coding-handbook/ch14\_evaluation/}]} \\
\noindent\rule{0.6\textwidth}{0.4pt}
\end{center}


\subsection{Attack 4: Denial of Wallet}

Craft requests that maximize token consumption per unit of legitimate value:

\begin{center}
\noindent\rule{0.6\textwidth}{0.4pt} \\
\vspace{0.3em}
\textit{[Code implementation is available in the companion coding handbook: \texttt{coding-handbook/ch14\_evaluation/}]} \\
\noindent\rule{0.6\textwidth}{0.4pt}
\end{center}

